\documentclass[10pt]{article}
\usepackage[T1]{fontenc}
\usepackage[utf8]{inputenc}
\usepackage[french]{babel}
\usepackage{amsmath,amssymb,mathrsfs,mathpazo}
\usepackage{geometry}
\geometry{a4paper,margin=2cm}
\usepackage{graphicx,array,tabularx,tabularray,colortbl,multirow,multicol,booktabs,enumitem,pifont,ulem,fancybox,diagbox,cancel}
\usepackage{pstricks,pstricks-add,pst-plot,pst-tree,pst-node,pst-text,pst-func,pst-math,pst-bspline,pst-3dplot,pst-eucl,pst-solides3d}
\let\pscalculate\undefined
\usepackage{pst-bezier}
\usepackage{tikz,pgfplots,tkz-tab}
\usetikzlibrary{arrows,arrows.meta,patterns,calc,positioning,trees,decorations.markings,shapes.geometric,shapes.misc,shapes.arrows,intersections,decorations.pathmorphing,decorations.pathreplacing,backgrounds,shadows,babel}
\pgfplotsset{compat=1.18}
\usepgfplotslibrary{fillbetween}
\usepackage[np]{numprint}
\usepackage{listings,makecell,pas-tableur,twoopt}
\let\barre\undefined
\usepackage[pstricks]{tablvar}
\definecolor{darkgreen}{RGB}{0,100,0}
\providecommand{\red}{\color{red}}
\providecommand{\blue}{\color{blue}}
\providecommand{\green}{\color{green}}
\providecommand{\black}{\color{black}}
\providecommand{\white}{\color{white}}
\providecommand{\gray}{\color{gray}}
\providecommand{\cyan}{\color{cyan}}
\providecommand{\magenta}{\color{magenta}}
\providecommand{\yellow}{\color{yellow}}
\providecommand{\violet}{\color{violet}}
\providecommand{\orange}{\color{orange}}
\providecommand{\euro}{\texteuro}
\providecommand{\EUR}{\texteuro}
\providecommand{\ieme}{\textsuperscript{e}}
\providecommand{\degres}{\ensuremath{^\circ}}
\providecommand{\degree}{\ensuremath{^\circ}}
\providecommand{\vv}[1]{\overrightarrow{#1}}
\providecommand{\Coord}[1]{#1}
\providecommand{\decosix}{}
\providecommand{\starredbullet}{\ensuremath{\bullet}}
\pagestyle{empty}
\setlength{\parindent}{0pt}
\usepackage{adjustbox}\DeclareUnicodeCharacter{2236}{:}\DeclareUnicodeCharacter{2190}{\ensuremath{\leftarrow}}
\providecommand{\R}{\mathbb{R}}
\providecommand{\N}{\mathbb{N}}
\providecommand{\Z}{\mathbb{Z}}
\providecommand{\Q}{\mathbb{Q}}
\providecommand{\D}{\mathbb{D}}
\providecommand{\C}{\mathbb{C}}
\providecommand{\e}{{\mathrm{e}}}
\providecommand{\vect}[1]{\overrightarrow{#1}}
\providecommand{\vectt}[1]{\overrightarrow{\mathrm{#1}}}
\providecommand{\vvt}[1]{\overrightarrow{\mathrm{#1}}}
\providecommand{\np}[2][]{{#2~#1}}
\providecommand{\vv}[1]{{\overrightarrow{#1}}}
\renewcommand{\d}{{\mathrm{d}}}
\providecommand{\degres}{{^\circ}}
\providecommand{\degre}{{^\circ}}
\providecommand{\blue}{\color{blue}}
\providecommand{\red}{\color{red}}
\providecommand{\enskip}{\quad}
\providecommand{\llbracket}{\lbrack\!\lbrack}
\providecommand{\rrbracket}{\rbrack\!\rbrack}
\providecommand{\euro}{{\text{€}}}
\providecommand{\no}{{\text{n°}}}
\providecommand{\interval}[2]{{[#1;#2]}}
\providecommand{\intervFO}[2]{{[#1;#2[}}
\providecommand{\dotfill}{\ldots}

\providecommand{\ds}{\displaystyle}
\providecommand{\textding}[1]{#1}
\providecommand{\Oij}{\(\mathrm{O};\vec{\imath},\vec{\jmath}\)}
\providecommand{\Oijk}{\(\mathrm{O};\vec{\imath},\vec{\jmath},\vec{k}\)}
\providecommand{\Ouv}{\(\mathrm{O};\vec{u},\vec{v}\)}
\providecommand{\intervFF}[2]{\ensuremath{[#1;#2]}}
\renewcommand{\euro}{\ifmmode\text{\texteuro}\else\texteuro\fi}
\DeclareUnicodeCharacter{00B0}{\ifmmode^{\circ}\else\textdegree\fi}
\DeclareUnicodeCharacter{20AC}{\ifmmode\text{\texteuro}\else\texteuro\fi}
% Énoncés : Éducation nationale. Transcription et corrigés : APMEP, Denis Vergès et contributeurs. Découpage et classement : CoopMaths, Sébastien Lozano, Éric Elter et contributeurs. Adaptation de la présentation : Maths Explorer.
% CC BY-SA : https://coopmaths.fr/www/about/
% https://coopmaths.fr/alea/static/bac/2022/tex/bac_2022_05_sujet1_madagascar_3.tex
\begin{document}
\noindent\textit{APMEP / CoopMaths. CC BY-SA. Présentation adaptée par Maths Explorer.}\par\medskip
\medskip

\begin{tabularx}{\linewidth}{|X|}\hline
\textbf{Principaux domaines abordés :} Suites;
Fonctions, Fonction exponentielle.\\ \hline
\end{tabularx}

\bigskip

\textbf{Partie A}

\medskip

On considère la fonction $f$ définie pour tout réel $x$ par:

\[f(x) = 1+x - \text{e}^{0,5x - 2}.\]

On admet que la fonction $f$ est dérivable sur $\R$. On note $f'$ sa dérivée.

\medskip

\begin{enumerate}
\item 
	\begin{enumerate}
		\item Déterminer la limite de la fonction $f$ en $- \infty$.
		\item Démontrer que, pour tout réel $x$ non nul, $f(x) = 1+ 0,5x\left(2 - \dfrac{\text{e}^{0,5x}}{0,5x} \times \text{e}^{-2}\right)$.
		
En déduire la limite de la fonction $f$ en $+\infty$.
	\end{enumerate}
\item 
	\begin{enumerate}
		\item Déterminer $f'(x)$ pour tout réel $x$.
		\item Démontrer que l'ensemble des solutions de l'inéquation $f'(x) < 0$ est l'intervalle

$]4 + 2\ln (2)~;~+\infty[$.
	\end{enumerate}	
\item Déduire des questions précédentes le tableau de variation de la fonction $f$ sur $\R$. On fera figurer la valeur exacte de l'image de $4 + 2\ln (2)$ par $f$.
\item Montrer que l'équation $f(x) = 0$ admet une unique solution sur l'intervalle $[-1~;~0]$.
\end{enumerate}

\bigskip

\textbf{Partie B}

\medskip

On considère la suite $\left(u_n\right)$ définie par $u_0 = 0$ et, pour tout entier naturel $n$,
\: $u_{n+1} = f\left(u_n\right)$ où $f$ est la fonction définie à la partie A.

\medskip

\begin{enumerate}
\item 
	\begin{enumerate}
		\item Démontrer par récurrence que, pour tout entier naturel $n$, on a : 
		
		\[u_n \leqslant u_{n+1} \leqslant 4.\]
		
		\item En déduire que la suite $\left(u_n\right)$ converge. On notera $\ell$ la limite.
	\end{enumerate}	
\item 
	\begin{enumerate}
		\item On rappelle que $f$ vérifie la relation $\ell = f(\ell)$.
		
Démontrer que $\ell  = 4$.
		\item ~
		
\begin{minipage}{0.48\linewidth}On considère la fonction \texttt{valeur} écrite ci-contre dans le langage Python :
		\end{minipage}\hfill 
\begin{minipage}{0.48\linewidth}
\begin{tabular}{|l|}\hline
\texttt{def valeur (a) :}\\
\quad \texttt{u = 0}\\
\quad \texttt{n = 0}\\
\quad \texttt{while u $\leqslant$ a:}\\
\qquad \texttt{u=1 + u - exp(0.5*u - 2)}\\
\qquad \texttt{n = n+1}\\ 
\quad return \texttt{n}\\ \hline
\end{tabular}
\end{minipage}

L'instruction \texttt{valeur}(3.99) renvoie la valeur 12.

Interpréter ce résultat dans le contexte de l'exercice.
	\end{enumerate}
\end{enumerate}

\bigskip

\end{document}
